\documentclass[../../../include/open-logic-section]{subfiles}

\begin{document}

\olfileid{sth}{spine}{rank}
\olsection{સ્તરાંક}

હવે આપણે તબક્કાઓને $V_\alpha$ તરીકે વ્યાખ્યાયિત કર્યા છે અને
દરેક ગણ કોઈક તબક્કાનો ઉપગણ છે તે જાણીએ છીએ. તેથી ગણનો
\emph{સ્તરાંક} વ્યાખ્યાયિત કરી શકીએ. સહજ રીતે કહીએ તો
$A$નો સ્તરાંક એ $A$ રચાય તેની પ્રથમ ક્ષણ છે. વધુ ચોક્કસ રીતે:

\begin{defn}\ollabel{defnsetrank}
દરેક ગણ $A$ માટે, $\setrank{A}$ એ એવી લઘુતમ ક્રમવાચક સંખ્યા
$\alpha$ છે કે $A
\subseteq V_\alpha$.
\end{defn}
\begin{prop}\ollabel{ranksexist}
	દરેક $A$ માટે $\setrank{A}$ અસ્તિત્વમાં છે.
\end{prop}
\begin{proof}
	સાબિતી સ્વાધ્યાય તરીકે છોડી છે.
\end{proof}
\begin{prob}
	\olref[sth][spine][rank]{ranksexist} સાબિત કરો.
\end{prob}
\noindent 
સ્તરાંકોના સુક્રમથી કેટલાંક મહત્ત્વનાં પરિણામો સાબિત કરી શકીએ:
\begin{prop}\ollabel{valphalowerrank}
કોઈ પણ ક્રમવાચક સંખ્યા $\alpha$ માટે,
$V_\alpha = \Setabs{x}{\setrank{x} \in \alpha}$.
\end{prop}

\begin{proof}
જો $\setrank{x} \in \alpha$ હોય, તો $x \subseteq V_{\setrank{x}} \in
V_\alpha$; અને $V_\alpha$ ઉપગણ-સંવૃત હોવાથી
(\olref[Valphabasic]{Valphabasicprops}નો વારંવાર ઉપયોગ કરીને)
$x \in V_\alpha$. ઊલટું, જો $x \in V_\alpha$ હોય, તો
$x \subseteq V_\alpha$, તેથી $\setrank{x} \leq \alpha$;
હવે સરળ અનંતાતીત અનુમાનથી $\setrank{x} \in \alpha$ મળે છે.
\sourcecorrection{OLMD-120}{મૂળની આ છેલ્લી પંક્તિમાં
``$x \notin V_\alpha$'' લખાયું છે, જે શરૂઆતની
``$x \in V_\alpha$'' ધારણાની વિરુદ્ધ જાય છે.
જોઈતો નિષ્કર્ષ $\setrank{x} \in \alpha$ છે: ખાલી
તબક્કે સાબિત કરવાનું કશું નથી; અનુગામી તબક્કે $x$
અગાઉના તબક્કાનો ઉપગણ હોવાથી તેનો સ્તરાંક અગાઉના
સૂચકાંકથી મોટો નથી; સીમા તબક્કે $x$ કોઈ પહેલાંના
તબક્કામાં આવે છે અને અનુમાનની ધારણા લાગુ પડે છે.}
\end{proof}
\begin{prob}
	\olref[sth][spine][rank]{valphalowerrank}માં ઉલ્લેખેલું
	સરળ અનંતાતીત અનુમાન પૂર્ણ કરો.
\end{prob}

\begin{prop}\ollabel{rankmemberslower}
જો $B \in A$ હોય, તો $\setrank{B} \in \setrank{A}$.
\end{prop}

\begin{proof}
\olref{valphalowerrank} મુજબ
$A \subseteq V_{\setrank{A}} = \Setabs{x}{\setrank{x} \in
\setrank{A}}$.
\end{proof}
\noindent
આ હકીકતથી એવું પરિણામ સ્થાપી શકીએ છીએ જેનાથી અનુમાનની
એક રીત દ્વારા \emph{બધા ગણો} વિશે વિધાનો સાબિત થાય છે:

\begin{thm}[$\in$-અનુમાનનું પ્રરૂપ] 
કોઈ પણ સૂત્ર $\phi$ માટે:
\[
	\forall A((\forall x \in A)\phi(x) \lif \phi(A)) \lif \forall A \phi(A).
\]
\end{thm}

\begin{proof}
આપણે પ્રતિવિધાન સાબિત કરીએ. ધારો કે $\lnot \forall A
\phi(A)$. અનંતાતીત અનુમાન
(\olref[ordinals][basic]{ordinductionschema}) મુજબ
લઘુતમ શક્ય સ્તરાંક ધરાવતું કોઈ બિન-$\phi$ મળે છે;
એટલે કે એવું $A$ છે કે $\lnot \phi(A)$ અને
$\forall x(\setrank{x} \in \setrank{A}
\lif \phi(x))$. હવે $x \in A$ હોય, તો
\olref{rankmemberslower} મુજબ $\setrank{x} \in
\setrank{A}$, તેથી $\phi(x)$; એટલે કે $(\forall x \in
A)\phi(x) \land \lnot \phi(A)$.
\end{proof}\noindent આ શક્તિશાળી પરિણામનો અનૌપચારિક અર્થ આ છે.
ગણનો દરેક !!{element} $\phi$ હોય ત્યારે ગણ પોતે પણ
$\phi$ હોય, તો અને માત્ર તો $\phi$ને \emph{વારસાગત}
કહીએ. ત્યારે $\in$-અનુમાન કહે છે: જો $\phi$ વારસાગત હોય,
તો દરેક ગણ $\phi$ છે.

સ્તરાંકોની ચર્ચા (હાલ માટે) પૂરી કરતાં પહેલાં, અગાઉ
કેટલીક વાર સૂચવેલાં થોડાં પરિણામો સાબિત કરીએ.

\begin{prop}\ollabel{ranksupstrict}
$\setrank{A} = \supstrict_{x \in A}\setrank{x}$.
\end{prop}

\begin{proof}
$\alpha = \supstrict_{x \in A}\setrank{x}$ લો.
\olref{rankmemberslower} મુજબ $\alpha \leq \setrank{A}$.
પરંતુ $x \in A$ હોય, તો $\setrank{x} \in \alpha$;
તેથી \olref{valphalowerrank} મુજબ $x \in V_\alpha$ અને આમ $A
\subseteq V_\alpha$, એટલે $\setrank{A} \leq \alpha$.
તેથી $\setrank{A} = \alpha$.
\end{proof}

\begin{cor}\ollabel{ordsetrankalpha}
કોઈ પણ ક્રમવાચક સંખ્યા $\alpha$ માટે $\setrank{\alpha} = \alpha$.
\end{cor}

\begin{proof}
અનંતાતીત અનુમાન માટે ધારો કે બધા $\beta \in \alpha$ માટે
$\setrank{\beta} = \beta$. હવે \olref{ranksupstrict} મુજબ
$\setrank{\alpha} = \supstrict_{\beta \in
\alpha}\setrank{\beta} = \supstrict_{\beta \in \alpha}\beta = \alpha$.
%	First note that $\setrank{\alpha} \neq \beta$ for any $\beta \in
%	\alpha$. For otherwise we would have $\beta = \setrank{\beta} \in
%	\setrank{\alpha} = \beta$ by \olref{rankmemberslower}, i.e.,
%	$\beta \in \beta$, a contradiction. 
%
%	Now note that $\alpha \subseteq V_\alpha$. For $\alpha =
%	\Setabs{\beta}{\beta \in \alpha}$ by
%	\olref[sfr][ordinals][basic]{ordissetofsmallerord}. And if
%	$\beta \in \alpha$ then $\beta \subseteq V_{\setrank{\beta}} =
%	V_\beta \in V_\alpha$, hence $\beta \in V_\alpha$, by
%	\olref[sfr][spine][Valphabasic]{Valphabasicprops} twice. 
\end{proof}

છેલ્લે, \olref[foundation]{sec}ના અંતે આપેલું વચન પૂર્ણ કરતી
ટૂંકી સાબિતી આપીએ: $\ZFminus$માં શરતી વિધાન
$\emph{Regularity} \Rightarrow \emph{Foundation}$ સાબિત થાય છે.
(ધ્યાન આપો કે આ સાબિતીમાં ``સ્તરાંક''નો ખ્યાલ અને
\olref{rankmemberslower} વાપરી શકીએ છીએ, કારણ કે આ વિભાગના
આરંભમાં કહ્યું તેમ તેમને $\ZFminus + \text{Regularity}$માં
રજૂ કરી શકાય છે.)

\begin{prop}[$\ZFminus + \text{Regularity}$માં કામ કરતાં]\ollabel{zfminusregularityfoundation} આધાર સાચો છે.
\end{prop}

\begin{proof}
$A \neq \emptyset$ લો અને લઘુતમ શક્ય સ્તરાંક ધરાવતું
કોઈ $B \in A$ પસંદ કરો. જો $c \in B$ હોય, તો
\olref{rankmemberslower} મુજબ $\setrank{c} \in \setrank{B}$;
તેથી $B$ની પસંદગી પ્રમાણે $c \notin A$.
\end{proof}

\end{document}
